\chapter{Garis sejajar}\label{chap:angle-sum}

{

\begin{wrapfigure}{o}{19mm}
\vskip-4mm
\centering
\includegraphics{mppics/pic-72}
\end{wrapfigure}

Ingat bahwa sebagai akibat Aksioma~\ref{def:birkhoff-axioms:1},
setiap dua garis berbeda $\ell$ dan $m$ memiliki tepat satu titik persekutuan atau tidak memiliki titik persekutuan sama sekali; lihat \ref{lem:line-line}.
Dalam kasus pertama, kedua garis tersebut \index{garis berpotongan}\emph{berpotongan} (secara ringkas $\ell\nparallel m$);
dalam kasus kedua, $\ell$ dan $m$ disebut \index{sejajar!garis}\emph{sejajar} (secara ringkas, \index{36@\hskip.4mm$\parallel$\hskip.4mm, \hskip.4mm$\nparallel$\hskip.4mm}$\ell\parallel m$);
selain itu, suatu garis selalu dianggap sejajar dengan dirinya sendiri.

}

Untuk menegaskan bahwa dua garis dalam suatu gambar sejajar, kita akan menandainya dengan tanda panah sejenis.



\begin{thm}[\abs]{Proposisi}\label{prop:perp-perp} Misalkan $\ell$, $m$, dan $n$ merupakan tiga garis.
Andaikan $n\perp \ell$ dan $n\perp m$.
Maka $\ell\parallel m$.
\end{thm}

\parit{Bukti.}
Andaikan sebaliknya;
yaitu, $\ell$ dan~$m$ berpotongan pada satu titik, sebutlah $Z$.
Maka menurut Teorema~\ref{perp:ex+un},
$\ell=m$.
Karena setiap garis sejajar dengan dirinya sendiri, kita memperoleh $\ell\parallel m$ --- suatu kontradiksi.
\qeds

\begin{thm}{Teorema}\label{thm:parallel}
Untuk setiap titik $P$ dan setiap garis $\ell$,
terdapat tepat satu garis $m$
yang melalui $P$ dan sejajar dengan~$\ell$.
\end{thm}

Teorema di atas memiliki dua bagian, yaitu keberadaan dan ketunggalan.
Dalam bukti ketunggalan, kita akan menggunakan metode segitiga sebangun.

\parit{Bukti; keberadaan.}
Terapkan Teorema~\ref{perp:ex+un} dua kali,
pertama untuk mengonstruksi garis $n$ melalui $P$ yang tegak lurus terhadap $\ell$,
dan kedua untuk mengonstruksi garis $m$ melalui $P$ yang tegak lurus terhadap~$n$.
Kemudian terapkan Proposisi~\ref{prop:perp-perp}.

\parit{Ketunggalan.}
Jika $P\in\ell$, maka $m=\ell$ menurut definisi garis sejajar.
Sekarang, andaikan $P\notin\ell$.

Mari kita mengonstruksi garis $n\ni P$ dan $m\ni P$ seperti dalam bukti keberadaan; jadi $n\perp \ell$, $m\perp n$, dan $m\parallel \ell$.

Andaikan terdapat garis lain $s\ni P$ yang sejajar dengan~$\ell$.
Pilih titik $Q\in s$ yang terletak pada sisi $m$ yang sama dengan~$\ell$.
Misalkan $R$ adalah titik kaki tegak lurus dari $Q$ pada~$n$.

\begin{figure}[!ht]
\centering
\includegraphics{mppics/pic-74}
\end{figure}

Misalkan $D$ adalah titik potong $n$ dan~$\ell$.
Menurut Proposisi~\ref{prop:perp-perp}, $(QR)\parallel m$.
Oleh karena itu, $Q$, $R$, dan $\ell$ terletak pada sisi yang sama terhadap~$m$.
Khususnya, $R\in [P D)$.

Pilih $Z\in [P Q)$ sedemikian sehingga
$$\frac{PZ}{PQ}=\frac{PD}{PR};$$
titik tersebut ada menurut Latihan~\ref{ex:trig==}.
Menurut kriteria kesebangunan SAS (atau, secara ekuivalen, menurut Aksioma~\ref{def:birkhoff-axioms:4}),
kita memperoleh $\triangle RPQ\sim \triangle DPZ$;
oleh karena itu, $(Z D)\perp(P D)$.
Akibatnya, $Z$ terletak pada $\ell$ dan $s$ --- suatu kontradiksi.\qeds

\begin{thm}{Akibat}\label{cor:parallel-1}
Andaikan $\ell$, $m$, dan $n$ merupakan garis-garis
sedemikian sehingga $\ell\parallel m$ dan $m\parallel n$.
Maka $\ell\parallel n$.
\end{thm}

\parit{Bukti.}
Andaikan sebaliknya; yaitu, $\ell\nparallel n$.
Maka terdapat titik $P\in \ell\cap n$.
Menurut Teorema~\ref{thm:parallel},
$n=\ell$ --- suatu kontradiksi.
\qeds

Menurut definisi, $\ell\parallel m$ jika dan hanya jika $m\z\parallel \ell$.
Oleh karena itu, menurut akibat di atas, ``$\parallel$'' merupakan suatu
\index{relasi ekuivalensi}\emph{relasi ekuivalensi}.
Artinya, untuk sembarang garis $\ell$, $m$, dan $n$, berlaku syarat-syarat berikut:
\begin{enumerate}[(i)]
\item $\ell\parallel \ell$;
\item jika $\ell\parallel m$, maka $m\parallel \ell$;
\item jika $\ell\parallel m$ dan $m\parallel n$, maka
$\ell\parallel n$.
\end{enumerate}

Latihan berikut memberikan kebalikan parsial Proposisi \ref{prop:perp-perp}.

\begin{thm}[!]{Latihan}\label{ex:perp-perp}
Misalkan $k$, $\ell$, $m$, dan $n$ merupakan garis-garis.
\begin{enumerate}[(a)]
 \item\label{ex:perp-perp:a}
 Andaikan $\ell\parallel m$ dan $m\perp n$.
 Tunjukkan bahwa $\ell\perp n$.
 \item\label{ex:perp-perp:b}
 Andaikan $k\perp \ell$, $\ell\parallel m$, dan $m\perp n$.
 Tunjukkan bahwa $k\parallel n$.
\end{enumerate}

\end{thm}

\begin{thm}{Latihan}\label{ex:construction-parallel}
Lakukan konstruksi dengan penggaris dan jangka untuk membuat garis yang melalui suatu titik yang diberikan dan sejajar dengan suatu garis yang diberikan.
\end{thm}

\section{Pencerminan terhadap titik}

Ingat bahwa jika $O$ merupakan titik tengah ruas $[XX']$,
maka kita mengatakan bahwa $X'$ merupakan \index{pencerminan terhadap titik}\emph{cerminan} $X$ terhadap titik $O$.
Selain itu, kita mengasumsikan bahwa $O'=O$; yaitu, $O$ merupakan cerminan dirinya sendiri terhadap dirinya sendiri.

Pernyataan berikut merupakan penyempurnaan Proposisi~\ref{prop:point-reflection}.

\begin{thm}[\abs]{Proposisi}\label{prop:point-reflection+}
Pencerminan terhadap suatu titik merupakan transformasi isometrik langsung.
\end{thm}

\parit{Bukti.}
Pilih suatu titik $O$.
Menurut Proposisi~\ref{prop:point-reflection}, pencerminan terhadap $O$ merupakan transformasi isometrik.
Perhatikan bahwa setiap sudut $\angle XOY$ bertolak belakang dengan cerminannya $\angle X'OY'$.
Menurut \ref{prop:vert}, $\measuredangle XOY\z=\measuredangle X'OY'$.
Oleh karena itu, pencerminan tersebut tidak mungkin merupakan transformasi isometrik tak langsung.

Menurut \ref{prop:direct-indirect}, setiap transformasi isometrik bersifat langsung atau tak langsung.
Oleh karena itu, pencerminan terhadap $O$ harus bersifat langsung.
\qeds

\begin{thm}{Latihan}\label{ex:reflections}
Andaikan $\angle AOB$ siku-siku.
Tunjukkan bahwa komposisi pencerminan terhadap garis $(OA)$ dan $(OB)$ merupakan pencerminan terhadap $O$.

Gunakan pernyataan ini dan Akibat~\ref{cor:reflection+angle} untuk menyusun bukti lain bagi~\ref{prop:point-reflection+}.
\end{thm}

{

\begin{wrapfigure}{r}{23mm}
\vskip-6mm
\centering
\includegraphics{mppics/pic-78}
\end{wrapfigure}


\begin{thm}{Teorema}\label{thm:parallel-point-reflection}
Misalkan $\ell$ dan $m$ merupakan garis-garis yang masing-masing melalui titik $P$ dan $Q$.
Andaikan $O$ merupakan titik tengah $[PQ]$.
Maka $m \parallel \ell$ jika dan hanya jika $m$ merupakan cerminan $\ell$ terhadap $O$.
\end{thm}

}

\parit{Bukti; bagian ``jika''.}
Andaikan $m$ merupakan cerminan $\ell$ terhadap $O$.
Andaikan $\ell\nparallel m$; yaitu, $\ell$ dan $m$ berpotongan pada satu titik $Z$.
Misalkan $Z'$ menyatakan cerminan $Z$ terhadap $O$.

\begin{figure}[!ht]
\centering
\includegraphics{mppics/pic-80}
\end{figure}

Titik $Z'$ terletak pada kedua garis $\ell$ dan $m$.
Akibatnya, $Z'=Z$, atau secara ekuivalen $Z=O$.
Dalam kasus ini, $O\in \ell$, sehingga cerminan $\ell$ terhadap $O$ adalah $\ell$ itu sendiri;
yaitu, $\ell=m$ dan khususnya $\ell\parallel m$ --- suatu kontradiksi.

\parit{Bagian ``hanya jika''.}
Misalkan $\ell'$ merupakan cerminan $\ell$ terhadap $O$.
Menurut bagian ``jika'' dari teorema, $\ell'\parallel \ell$.
Kedua garis $\ell'$ dan $m$ melalui $Q$.
Menurut ketunggalan garis sejajar (\ref{thm:parallel}), jika $m\parallel \ell$, maka $\ell'=m$; dengan demikian, pernyataan terbukti.
\qeds

% Reflow id-ID: the inherited forced page break is omitted; content and order are unchanged.

\section[Sifat transversal]{\hspace{1em}Sifat transversal}

{

\begin{wrapfigure}{r}{25mm}
\centering
\vskip-10mm
\includegraphics{mppics/pic-82}
\end{wrapfigure}

Jika garis $t$ memotong masing-masing garis $\ell$ dan $m$ pada satu titik, maka kita menyebut $t$ sebagai \index{transversal}\emph{transversal} bagi $\ell$ dan~$m$.
Sebagai contoh, pada gambar, garis $(CB)$ merupakan transversal
bagi $(AB)$ dan~$(CD)$.

}

\begin{thm}{Sifat transversal}\label{thm:parallel-2}
Misalkan $A$, $B$, $C$, dan $D$ merupakan titik-titik sedemikian sehingga $A\ne B$, $B\ne C$, dan $C\ne D$.
Maka $(AB)\parallel(C D)$ jika dan hanya jika
$$2\cdot(\measuredangle A B C+\measuredangle B C D)\equiv 0.
\eqlbl{A B C + B C D}$$
Secara ekuivalen,
$$\measuredangle A B C+\measuredangle B C D
\equiv
0
\quad
\text{atau}
\quad
\measuredangle A B C+\measuredangle B C D
\equiv
\pi.$$

Selain itu, jika $(AB)\ne(C D)$, maka dalam kasus pertama,
$A$ dan $D$ terletak pada sisi-sisi yang berlawanan terhadap $(BC)$;
dalam kasus kedua,
$A$ dan $D$ terletak pada sisi yang sama terhadap~$(BC)$.
\end{thm}



\parit{Bukti; bagian ``hanya jika''.}
Misalkan $O$ menyatakan titik tengah $[BC]$.

Andaikan $(AB)\parallel(C D)$.
Menurut Teorema~\ref{thm:parallel-point-reflection},
$(CD)$ merupakan cerminan $(AB)$ terhadap $O$.

\begin{wrapfigure}{r}{31mm}
\vskip-4mm
\centering
\includegraphics{mppics/pic-84}
\end{wrapfigure}

Misalkan $A'$ merupakan cerminan $A$ terhadap $O$.
Maka $A'\in (CD)$ dan menurut Proposisi~\ref{prop:point-reflection+}, kita memperoleh
\[\measuredangle ABO=\measuredangle A'CO.\eqlbl{A B O = A' C O}\]
Perhatikan bahwa
\[\measuredangle ABO\equiv\measuredangle ABC,
\qquad
\measuredangle A'CO\equiv-\measuredangle BCA'.\eqlbl{eq:A B O= A B C}
\]
Karena $A'$, $C$, dan $D$ terletak pada satu garis, Latihan~\ref{ex:ABCO-line} menyiratkan bahwa
\[2\cdot \measuredangle BCD\equiv 2\cdot \measuredangle BCA'.\eqlbl{eq:2BCD= 2BCA'}\]
Akhirnya, \ref{A B O = A' C O}, \ref{eq:A B O= A B C}, dan \ref{eq:2BCD= 2BCA'} menyiratkan \ref{A B C + B C D}.
\qeds

\parit{Bagian ``jika''.}
Menurut Teorema~\ref{thm:parallel}, terdapat tepat satu garis $(CD)$ melalui $C$ yang sejajar dengan $(AB)$.
Dari bagian ``hanya jika'', kita mengetahui bahwa \ref{A B C + B C D} berlaku.

Di sisi lain, terdapat \textit{tepat satu} garis $(CD)$ sedemikian sehingga \ref{A B C + B C D} berlaku.
Memang, andaikan terdapat dua garis semacam itu, $(CD)$ dan $(CD')$, maka
$$2\cdot(\measuredangle A B C+\measuredangle B C D)\equiv 2\cdot(\measuredangle A B C+\measuredangle B C D')\equiv0.
$$
Oleh karena itu,
$2\cdot\measuredangle B C D\equiv 2\cdot\measuredangle B C D'$
dan menurut Latihan~\ref{ex:ABCO-line}, $D'\in (CD)$, atau secara ekuivalen garis $(CD)$ berimpit dengan $(CD')$.

Oleh karena itu, jika \ref{A B C + B C D} berlaku, maka $(CD)\parallel (AB)$.


\parit{Pernyataan terakhir.}
Jika $(AB)\ne(C D)$ serta $A$ dan $D$ terletak pada sisi-sisi yang berlawanan terhadap $(BC)$, maka $\angle ABC$ dan $\angle BCD$ memiliki tanda yang berlawanan.
Oleh karena itu,
\[-\pi\z<\measuredangle A B C+\measuredangle B C D<\pi.\]
Dengan menerapkan \ref{A B C + B C D}, kita memperoleh $\measuredangle A B C+\measuredangle B C D=0$.

Demikian pula, jika $A$ dan $D$ terletak pada sisi yang sama terhadap $(BC)$,
maka $\angle ABC$ dan $\angle BCD$ memiliki tanda yang sama.
Oleh karena itu,
\[0<|\measuredangle A B C+\measuredangle B C D|<2\cdot\pi,\]
dan \ref{A B C + B C D} menyiratkan bahwa $\measuredangle A B C+\measuredangle B C D\z\equiv\pi$.
\qeds

\begin{thm}{Latihan}\label{ex:parallel-angles}
Andaikan $(AB) \parallel (A'B')$ dan $(BC) \parallel (B'C')$.
Tunjukkan bahwa
$2\cdot \measuredangle ABC \equiv 2\cdot \measuredangle A'B'C'$.
Simpulkan bahwa salah satu dari
\[
\measuredangle ABC \equiv \measuredangle A'B'C' \quad \text{atau} \quad \measuredangle ABC \equiv \pi + \measuredangle A'B'C',
\]
berlaku, dan gambarlah sepasang sudut untuk mengilustrasikan setiap kasus.
\end{thm}


\begin{thm}[!]{Latihan}\label{ex:smililar+parallel}
Misalkan $\triangle ABC$ merupakan segitiga tak degenerat, dan andaikan $P$ terletak di antara $A$ dan $B$.
Andaikan suatu garis $\ell$ melalui $P$ dan sejajar dengan $(AC)$.
Tunjukkan bahwa $\ell$ memotong sisi $[BC]$ pada satu titik lain, sebutlah $Q$, dan
\[\triangle ABC\sim\triangle PBQ.\]
Khususnya, $\tfrac{PB}{AB}=\tfrac{QB}{CB}$.
\end{thm}

\begin{thm}{Latihan}\label{ex:trisection}
Bagilah suatu ruas yang diberikan menjadi tiga bagian sama panjang dengan penggaris dan jangka.
\end{thm}

\section{Sudut-sudut segitiga}

\begin{thm}{Teorema}\label{thm:3sum}
Pada setiap $\triangle A B C$, berlaku
$$\measuredangle A B C+ \measuredangle B C A + \measuredangle C A B \equiv \pi.$$

\end{thm}

\parit{Bukti.}
Jika $\triangle A B C$ degenerat, maka kesamaan tersebut mengikuti Akibat~\ref{cor:degenerate=pi}.
Sekarang, andaikan $\triangle A B C$ tak degenerat.

Misalkan $X$ merupakan cerminan $C$ terhadap titik tengah $M$ dari $[AB]$.
Menurut Proposisi~\ref{prop:point-reflection+},
$\measuredangle BAX\z=\measuredangle ABC$.
Perhatikan bahwa $(AX)$ merupakan cerminan $(CB)$ terhadap $M$;
oleh karena itu, menurut Teorema~\ref{thm:parallel-point-reflection}, $(AX)\z\parallel (CB)$.


\begin{wrapfigure}{o}{23mm}
\centering
\vskip-0mm
\includegraphics{mppics/pic-86}
\end{wrapfigure}

Karena $[BM]$ dan $[MX]$ tidak memotong $(CA)$,
titik-titik $B$, $M$, dan $X$ terletak pada sisi yang sama terhadap $(CA)$.
Dengan menerapkan sifat transversal pada transversal $(CA)$ bagi $(AX)$ dan $(CB)$, kita memperoleh
\[\measuredangle BCA+\measuredangle CAX\equiv \pi.\eqlbl{eq:ABC+CAB}\]

Karena $\measuredangle BAX=\measuredangle ABC$,
berlaku
\[\measuredangle CAX\equiv\measuredangle CAB+\measuredangle ABC.\]
Identitas terakhir ini dan \ref{eq:ABC+CAB} menyiratkan teorema.\qeds

\begin{thm}[!]{Latihan kelas}\label{ex:|3sum|}
Tunjukkan bahwa
$$|\measuredangle A B C|+ |\measuredangle B C A| + |\measuredangle C A B| = \pi$$
untuk setiap $\triangle ABC$.
\end{thm}

{

\begin{wrapfigure}{r}{30mm}
\vskip-0mm
\centering
\includegraphics{mppics/pic-88}
\vskip4mm
\includegraphics{mppics/pic-90}
%\vskip4mm
%\includegraphics{mppics/pic-92}
\end{wrapfigure}

\begin{thm}{Latihan}\label{ex:pent}
Misalkan $\triangle ABC$ merupakan segitiga tak degenerat.
Andaikan terdapat titik $D\in [BC]$
sedemikian sehingga
\[\measuredangle BAD=\measuredangle DAC,
\quad
BA=AD=DC.\]
Tunjukkan bahwa $\measuredangle CBA= \measuredangle BAC$ dan $AC=BC$.
\end{thm}


\begin{thm}{Latihan}\label{ex:right-isos}
Misalkan $\triangle ABC$ merupakan segitiga sama kaki tak degenerat dengan alas~$[AC]$.
Andaikan $D$ merupakan cerminan $A$ terhadap $B$.
Tunjukkan bahwa $\angle ACD$ siku-siku.
\end{thm}



%\begin{thm}{Exercise}\label{ex:pi/4-isos}
%Let $\triangle ABC$ be an isosceles nondegenerate triangle with base~$[AC]$.
%Assume that a circle is passing thru $A$, centered at a point on $[AB]$, and tangent to $(BC)$ at the point~$X$.
%Show that $\measuredangle CAX\z=\pm\tfrac\pi4$.
%\end{thm}

}

\begin{thm}[!]{Latihan}\label{ex:quadrangle}
Tunjukkan bahwa untuk setiap segiempat $ABCD$, berlaku
$$\measuredangle ABC+\measuredangle BCD+\measuredangle CDA+\measuredangle DAB\equiv 0.$$

\end{thm}


\section{Jajar genjang}

{

\begin{wrapfigure}{r}{26mm}
\vskip-10mm
\centering
\includegraphics{mppics/pic-94}
\end{wrapfigure}

Suatu segiempat $ABCD$ di bidang Euklides disebut \index{segiempat!segiempat tak degenerat}\index{tak degenerat!segiempat}\emph{tak degenerat} jika tidak ada tiga dari titik-titik $A,B,C,D$ yang terletak pada satu garis.

}

Suatu segiempat tak degenerat disebut \index{jajar genjang}\emph{jajar genjang}
jika sisi-sisi yang berhadapan saling sejajar.



\begin{thm}{Lema}\label{lem:parallelogram}
Setiap jajar genjang \index{simetri pusat}\emph{simetris pusat} terhadap titik tengah salah satu diagonalnya;
yaitu, pencerminan terhadap titik tengah tersebut memetakan jajar genjang itu ke dirinya sendiri.

Khususnya, jika $\square A B C D$ merupakan jajar genjang, maka
\begin{enumerate}[(a)]
\item\label{lem:parallelogram:midpoint} diagonal-diagonalnya $[AC]$ dan $[BD]$ saling berpotongan di titik tengah keduanya;
\item $\measuredangle A B C= \measuredangle C D A$;
\item $AB=CD$.
\end{enumerate}
\end{thm}

{

\begin{wrapfigure}{r}{33mm}
\centering
\includegraphics{mppics/pic-96}
\end{wrapfigure}

\parit{Bukti.} Misalkan $\square A B C D$ merupakan jajar genjang.
Misalkan $M$ menyatakan titik tengah $[AC]$.

Karena $(AB)\parallel (CD)$, Teorema~\ref{thm:parallel-point-reflection} menyiratkan bahwa $(CD)$ merupakan cerminan $(AB)$ terhadap $M$.
Dengan cara yang sama, $(BC)$ merupakan cerminan $(DA)$ terhadap $M$.
Karena $\square A B C D$ tak degenerat, diperoleh bahwa $D$ merupakan cerminan $B$ terhadap $M$; dengan kata lain, $M$ merupakan titik tengah $[BD]$.

Pernyataan-pernyataan lainnya mengikuti karena pencerminan terhadap $M$ merupakan transformasi isometrik langsung pada bidang (lihat \ref{prop:point-reflection+}).
\qeds

}

\begin{thm}{Latihan}\label{ex:4parallels}
Andaikan titik $P$ berbeda dari titik-titik sudut suatu jajar genjang $ABCD$.
Misalkan $a$, $b$, $c$, dan $d$ merupakan garis-garis yang melalui titik $A$, $B$, $C$, dan $D$ sedemikian sehingga
$a\parallel (CP)$,
$b\parallel (DP)$,
$c\parallel (AP)$, dan
$d\parallel (BP)$.
Tunjukkan bahwa garis-garis $a$, $b$, $c$, dan $d$ berpotongan pada satu titik.
\end{thm}


\begin{thm}{Latihan}\label{ex:romb}
Andaikan $ABCD$ merupakan segiempat sedemikian sehingga
\[AB=CD=BC=DA.\]
Tunjukkan bahwa $ABCD$ merupakan jajar genjang.
\end{thm}

Segiempat seperti dalam latihan di atas disebut \index{belah ketupat}\emph{belah ketupat}.

Suatu segiempat $ABCD$ disebut \index{persegi panjang}\emph{persegi panjang} jika sudut-sudut $ABC$, $BCD$, $CDA$, dan $DAB$ siku-siku.
Menurut sifat transversal (\ref{thm:parallel-2}),
setiap persegi panjang merupakan jajar genjang.

Persegi panjang yang semua sisinya sama panjang disebut \index{persegi}\emph{persegi}.

\begin{thm}{Latihan}\label{ex:rectangle}
Tunjukkan bahwa suatu jajar genjang $ABCD$ merupakan persegi panjang
jika dan hanya jika $AC=BD$.
\end{thm}

\begin{thm}{Latihan}\label{ex:romb2}
Tunjukkan bahwa suatu jajar genjang $ABCD$ merupakan belah ketupat
jika dan hanya jika $(AC)\perp (BD)$.
\end{thm}

Andaikan $\ell\parallel m$ dan $X,Y\in m$.
Misalkan $X'$ dan $Y'$ berturut-turut menyatakan titik kaki tegak lurus dari $X$ dan $Y$ pada~$\ell$.
Perhatikan bahwa $\square XYY'X'$ merupakan persegi panjang.
Menurut Lema~\ref{lem:parallelogram}, $XX'=YY'$.
Artinya, setiap titik pada $m$ memiliki jarak yang sama dari $\ell$.
Jarak ini disebut \index{jarak!antara garis sejajar}\emph{jarak antara} $\ell$ dan~$m$.

\begin{thm}{Latihan}\label{ex:inscribed-rhombus}
Suatu belah ketupat $AXYZ$ dengan panjang sisi $s$ terinskripsi dalam segitiga tak degenerat $ABC$ sedemikian sehingga $X\in [AB]$, $Y\in [BC]$, dan $Z\in[CA]$.
Tunjukkan bahwa $s=\sqrt{XB\cdot ZC}$.
\end{thm}

\section{Metode koordinat}

Latihan berikut penting;
latihan tersebut menunjukkan bahwa definisi aksiomatik kita sesuai dengan model yang diuraikan dalam Bagian~\ref{page:model}.


\begin{thm}{Latihan}\label{ex:coordinates}
Misalkan $\ell$ dan $m$ merupakan dua garis yang saling tegak lurus di bidang Euklides.
Untuk suatu titik $P$, misalkan $P_\ell$ dan $P_m$ berturut-turut menyatakan titik kaki tegak lurus dari $P$ pada $\ell$ dan $m$.

\begin{enumerate}[(a)]
\item Tunjukkan bahwa untuk setiap $X\in \ell$ dan $Y\in m$, terdapat tepat satu titik $P$ sedemikian sehingga $P_\ell=X$ dan $P_m=Y$.
\end{enumerate}

\begin{enumerate}[(a)]\addtocounter{enumi}{1}
\item
Tunjukkan bahwa
$PQ^2=P_\ell Q_\ell^2+P_mQ_m^2$
untuk setiap pasangan titik $P$ dan~$Q$.
\end{enumerate}

\begin{enumerate}[(a)]\addtocounter{enumi}{2}
\item Simpulkan bahwa bidang tersebut isometrik dengan $(\mathbb{R}^2,d_2)$; lihat \ref{sec:Metric spaces}.
\end{enumerate}

\end{thm}

\begin{wrapfigure}{r}{37mm}
\centering
\includegraphics{mppics/pic-98}
\end{wrapfigure}

Setelah latihan ini diselesaikan, kita dapat menerapkan
metode koordinat
untuk menyelesaikan sembarang masalah dalam geometri bidang Euklides.
Metode ini kuat dan universal;
metode tersebut akan dikembangkan lebih lanjut dalam Bab~\ref{chap:complex}.

\begin{thm}{Latihan}\label{ex:abc}
Gunakan Latihan~\ref{ex:coordinates}
untuk memberikan bukti alternatif bagi Teorema~\ref{thm:abc} di bidang Euklides.

Yaitu, buktikan bahwa untuk bilangan-bilangan real $a$, $b$, dan $c$ yang memenuhi
 $$0<a\le b\le c\le a+b,$$
terdapat suatu segitiga $ABC$
sedemikian sehingga $a=BC$, $b=CA$, dan $c=AB$.
\end{thm}

\begin{thm}[!]{Latihan}\label{ex:line-coord}
Perhatikan dua titik berbeda $A=(x_A,y_A)$ dan $B\z=(x_B,y_B)$ pada bidang koordinat.
Tunjukkan bahwa garis sumbu $[AB]$ dinyatakan oleh persamaan
\[2\cdot (x_B-x_A)\cdot x+2\cdot (y_B-y_A)\cdot y=x_B^2+y_B^2-x_A^2-y_A^2.\]

Simpulkan bahwa suatu garis dapat didefinisikan sebagai himpunan bagian bidang koordinat dalam salah satu bentuk berikut:
\begin{enumerate}[(a)]
\item Himpunan solusi suatu persamaan $a\cdot x+b\cdot y=c$
untuk konstanta $a$, $b$, dan $c$ sedemikian sehingga $a\ne 0$ atau $b\ne0$.
\item\label{ex:line-coord:parameter} Himpunan titik $(a\cdot t+c,b\cdot t+d)$ untuk semua $t\in \mathbb{R}$, dengan $a$, $b$, $c$, dan $d$ merupakan konstanta yang memenuhi $a\ne 0$ atau $b\ne0$.
\end{enumerate}

\end{thm}

\section{Lingkaran Apollonius}\label{sec:Apollonian circle}
Latihan-latihan dalam bagian ini mengilustrasikan metode koordinat; latihan-latihan tersebut tidak akan digunakan lagi selanjutnya.

\begin{thm}{Latihan}\label{ex:circle-coord}
Tunjukkan bahwa untuk nilai-nilai real tetap $a$, $b$, dan $c$, persamaan
\[x^2+y^2+a\cdot x+b\cdot y+c=0\]
menggambarkan suatu lingkaran, suatu titik, atau himpunan kosong.

Tunjukkan bahwa jika persamaan tersebut menggambarkan suatu lingkaran, maka $(-\tfrac a2,-\tfrac b2)$ merupakan pusatnya,
dan $r\z=\tfrac12\cdot \sqrt{a^2+b^2-4\cdot c}$ merupakan jari-jarinya.
\end{thm}

\begin{thm}{Latihan}\label{ex:apolonnius}
Gunakan latihan sebelumnya untuk menunjukkan bahwa, untuk suatu bilangan real positif $k\ne1$,
lokus titik-titik $M$ yang memenuhi $AM=k\cdot BM$,
dengan $A$ dan $B$ dua titik berbeda,
merupakan suatu lingkaran.
\end{thm}

\begin{figure}[!ht]
\centering
\includegraphics{mppics/pic-100}
\end{figure}

Lingkaran dalam latihan di atas merupakan contoh dari apa yang disebut \index{lingkaran Apollonius}\emph{lingkaran Apollonius dengan fokus $A$ dan $B$}.
Beberapa lingkaran semacam ini untuk nilai $k$ yang berbeda ditunjukkan pada gambar;
untuk $k=1$, lingkaran tersebut menjadi garis sumbu $[AB]$.


\begin{thm}{Latihan}\label{ex:apolonnius-construction}
Andaikan $A\ne B$, $M\notin\{A,B\}$, dan $AM\ne BM$.
Dengan penggaris dan jangka, konstruksilah lingkaran Apollonius dengan fokus $A$ dan $B$ yang diberikan dan melalui suatu titik $M$ yang diberikan.
\end{thm}
